\chapter{Three Classic Classifiers: $k$-NN, Naive Bayes and Fisher's LDA}
\companion{StatQuest: Naive Bayes, Clearly Explained}{\ytid{O2L2Uv9pdDA}}

The InfoEdge hire was asked to build Naive Bayes up from Bayes' theorem ``with intuition, not formulas'', and to explain
Fisher's linear discriminant (``push class means apart, shrink the variance within each class''). $k$-nearest neighbours
is the simplest possible learner and a standard warm-up question. Each is a different philosophy: $k$-NN memorises,
Naive Bayes models how data are generated, LDA finds the best projection.

\section{$k$-nearest neighbours}

\subsection{The idea}
To classify a new point, find the $k$ training points closest to it and take a majority vote (or average their targets
for regression). No training at all: the model \emph{is} the data. This is called a \textbf{lazy} or
\textbf{instance-based} learner.

\begin{example}[: 3-NN by hand]
Training points (experience, skills matched) with labels ``shortlisted'' (S) or not (N): $A(1,2)$ N, $B(2,3)$ N,
$C(5,8)$ S, $D(6,7)$ S, $E(4,6)$ S. New candidate $Q = (4, 5)$. Euclidean distances: $A$: $\sqrt{9+9} = 4.24$; $B$:
$\sqrt{4+4} = 2.83$; $C$: $\sqrt{1+9} = 3.16$; $D$: $\sqrt{4+4} = 2.83$; $E$: $1$. Nearest three: $E$ (S), $B$ (N), $D$ (S).
Vote: S wins, 2 to 1.
\end{example}

\subsection{Choosing $k$}
\begin{itemize}
\item $k = 1$: the boundary follows every training point; zero training error; high variance (one mislabelled point
  creates an island).
\item Large $k$: smooth boundary; high bias; with $k = n$ it always predicts the majority class.
\item Choose $k$ by cross-validation; odd $k$ avoids ties in binary problems; distance-weighted voting ($1/d$) gives
  closer neighbours more say.
\end{itemize}

\subsection{Distances and scaling}
Euclidean $\sqrt{\sum(x_j - z_j)^2}$, Manhattan $\sum|x_j - z_j|$, cosine (for text and embeddings), Hamming (for
categorical). \textbf{Scaling is essential}: a feature measured in rupees will dominate one measured in years.

\subsection{The curse of dimensionality}
In high dimensions almost all points are far from each other and distances become similar: the nearest neighbour is barely
nearer than the farthest. Example: to capture 10\% of uniformly spread data in 10 dimensions, a cube around the query must span
$0.1^{1/10} = 79\%$ of each axis: ``local'' is no longer local. $k$-NN therefore needs dimensionality reduction or learned
embeddings to work in high dimensions.

\subsection{Cost}
Training $O(1)$; prediction $O(nd)$ per query by brute force. Speed-ups: KD-trees and ball trees (good up to about 20
dimensions), and \textbf{approximate nearest neighbour} indexes (HNSW, IVF, product quantisation; libraries FAISS, ScaNN)
for millions of embeddings. The PremiumX slide's ``vector index'' is exactly this.

\section{Naive Bayes, built from Bayes' theorem}

\subsection{Bayes' theorem in words}
We want $P(\text{class}\mid\text{evidence})$. Bayes' theorem turns it around:
\[ P(C\mid x) = \frac{P(x\mid C)\,P(C)}{P(x)}. \]
\begin{itemize}
\item $P(C)$, the \textbf{prior}: how common the class is before looking at the evidence (2\% of messages are spam).
\item $P(x\mid C)$, the \textbf{likelihood}: how typical this evidence is for that class (``free offer'' appears in 30\% of
  spam and 1\% of normal mail).
\item $P(x)$, the \textbf{evidence}: the same for every class, so for choosing a class we can ignore it and compare
  $P(x\mid C)P(C)$ across classes.
\item $P(C\mid x)$, the \textbf{posterior}: the updated belief.
\end{itemize}

\subsection{The problem: too many combinations}
The evidence is many features at once: thousands of words. Estimating $P(x_1, x_2, \dots, x_d\mid C)$ directly would need
data for every combination of words, impossible.

\subsection{The naive assumption}
Assume the features are \textbf{conditionally independent given the class}: once you know it is spam, knowing it contains
``free'' tells you nothing extra about whether it contains ``offer''. Then the joint likelihood factorises:
\[ P(x_1,\dots,x_d\mid C) = \prod_{j=1}^dP(x_j\mid C), \qquad \hat C = \argmax_C\ P(C)\prod_jP(x_j\mid C). \]
Now each $P(x_j\mid C)$ is a simple count. The assumption is almost always false (``free'' and ``offer'' co-occur), but the
\emph{ranking} of classes is often still right, which is all classification needs.

\begin{example}[: spam by hand]
Prior: $P(S) = 0.3$, $P(H) = 0.7$. Word probabilities: $P(\text{free}\mid S) = 0.4$, $P(\text{free}\mid H) = 0.05$;
$P(\text{meeting}\mid S) = 0.05$, $P(\text{meeting}\mid H) = 0.3$. Message: ``free meeting''.
Spam score: $0.3\times0.4\times0.05 = 0.006$. Ham score: $0.7\times0.05\times0.3 = 0.0105$.
Normalise: $P(S\mid\text{msg}) = 0.006/0.0165 = 0.36$. Classified as ham. ``Meeting'' outweighed ``free''.
\end{example}

\subsection{Variants}
\begin{itemize}
\item \textbf{Multinomial NB}: features are word counts; $P(w\mid C) = \frac{\text{count}(w, C) + \alpha}{\text{total words in }C + \alpha|V|}$.
\item \textbf{Bernoulli NB}: features are word presence/absence; also penalises the absence of words typical of a class.
\item \textbf{Gaussian NB}: numeric features modelled as a normal distribution per class.
\item \textbf{Complement NB}: estimates from the complement of each class; better with imbalanced text.
\end{itemize}

\subsection{Laplace smoothing: the zero-frequency problem}
If ``lottery'' never appeared in ham training mail, $P(\text{lottery}\mid H) = 0$, and one occurrence makes the whole ham
product 0, whatever else the message says. Add $\alpha$ (usually 1) to every count:
\begin{example}[: smoothing]
Spam class has 20 word tokens in total, vocabulary size 10, ``offer'' appears 3 times in spam.
Unsmoothed: $3/20 = 0.15$. Laplace ($\alpha = 1$): $(3+1)/(20+10) = 4/30 = 0.133$. A word never seen in spam:
$1/30 = 0.033$ instead of 0.
\end{example}
Smoothing is a MAP estimate with a Dirichlet (uniform-ish) prior (Chapter 12).

\subsection{Log space}
Multiplying hundreds of probabilities like 0.001 underflows to 0 in floating point. Use sums of logs:
$\log P(C) + \sum_j\log P(x_j\mid C)$. The argmax is unchanged because $\log$ is increasing.

\subsection{Why it works surprisingly well}
Very few parameters (one per word per class), so low variance: excellent with little data and huge vocabularies; trains in one
pass; predicts fast. Its probabilities are poorly calibrated (the independence assumption double-counts correlated evidence, so
outputs are pushed toward 0 and 1), but its rankings are often good.

\section{Fisher's Linear Discriminant Analysis}

\subsection{The goal}
Project data onto a single direction $\mathbf w$ (a line) such that, along that line, the two classes are as separated as
possible. ``Separated'' means two things at once: the class \textbf{means} far apart, and each class \textbf{tightly bunched}.

\begin{intuition}[: why means alone are not enough]
Imagine two cigar-shaped clouds lying side by side, parallel. The direction joining their means may run along the cigars'
length, where each cloud is very spread out, so the projections overlap heavily. A slightly rotated direction, across the
cigars, separates them cleanly even though the means are a bit closer along it. Fisher's criterion finds that rotation.
\end{intuition}

\subsection{The criterion}
For a direction $\mathbf w$, the projected class means are $\tilde m_k = \mathbf w^\top\mathbf m_k$ and the projected within-class
scatters are $\tilde s_k^2 = \sum_{i\in k}(\mathbf w^\top\mathbf x_i - \tilde m_k)^2$. Fisher maximises
\[ J(\mathbf w) = \frac{(\tilde m_1 - \tilde m_2)^2}{\tilde s_1^2 + \tilde s_2^2} = \frac{\mathbf w^\top S_B\mathbf w}{\mathbf w^\top S_W\mathbf w}, \]
with between-class scatter $S_B = (\mathbf m_1 - \mathbf m_2)(\mathbf m_1 - \mathbf m_2)^\top$ and within-class scatter
$S_W = \sum_k\sum_{i\in k}(\mathbf x_i - \mathbf m_k)(\mathbf x_i - \mathbf m_k)^\top$. Setting the derivative to zero gives the solution
\[ \boxed{\mathbf w\propto S_W^{-1}(\mathbf m_1 - \mathbf m_2).} \]
In words: start from the direction joining the means, then ``undo'' the shape of the clouds by multiplying by the inverse of the
within-class scatter.

\begin{example}[: Fisher direction in 2-D]
Class means $\mathbf m_1 = (2, 2)$, $\mathbf m_2 = (4, 3)$, so $\mathbf m_1 - \mathbf m_2 = (-2, -1)$. Suppose
$S_W = \begin{pmatrix}4 & 0\\0 & 1\end{pmatrix}$ (each class spread out along $x_1$). $S_W^{-1} = \text{diag}(0.25, 1)$.
$\mathbf w\propto(-0.5, -1)$. Joining the means suggests weighting $x_1$ twice as much as $x_2$; Fisher weights $x_2$ twice as much as
$x_1$, because along $x_1$ the classes are noisy.
\end{example}

\subsection{More than two classes, and LDA as a classifier}
With $K$ classes, LDA finds up to $K - 1$ directions: the top eigenvectors of $S_W^{-1}S_B$. It is a \textbf{supervised}
dimensionality reduction, in contrast to PCA (unsupervised, maximises total variance, ignores labels).
LDA is also a probabilistic classifier: assume each class is Gaussian with a \emph{shared} covariance $\Sigma$; then the Bayes
decision boundary is linear, and its direction is exactly Fisher's $\Sigma^{-1}(\mathbf m_1 - \mathbf m_2)$. With class-specific
covariances the boundary is quadratic: \textbf{QDA}.

\subsection{Limits}
Linear boundaries only; assumes roughly Gaussian classes with similar covariance for optimality; $S_W$ must be invertible
(regularise with shrinkage when $d$ is large); at most $K-1$ dimensions.

\begin{practical}[: these three at InfoEdge]
\begin{itemize}
\item \textbf{$k$-NN / vector search}: ``similar jobs'', ``similar candidates'' and the dense half of PremiumX's hybrid retrieval
  are nearest-neighbour search over embeddings with an ANN index.
\item \textbf{Naive Bayes}: a strong, instant baseline for classifying job titles, routing support tickets, or spam/abuse detection in
  Jeevansathi messages.
\item \textbf{LDA}: a quick supervised projection to visualise how separable ``placed'' vs ``not placed'' students are, and a
  well-behaved linear classifier for small datasets.
\end{itemize}
\end{practical}

\stuck{StatQuest: K-nearest neighbors, Clearly Explained}{\ytid{HVXime0nQeI}}
\section{Interview questions}

\begin{qa}{\pyq{INT: Naive Bayes from Bayes' theorem} Build Naive Bayes from Bayes' theorem. What assumption does it make, why is it
needed, and how is it used in practice?}
\oneline{Naive Bayes applies Bayes' theorem to get $P(\text{class}\mid\text{features})\propto P(\text{class})\prod_jP(x_j\mid\text{class})$,
which requires assuming the features are independent given the class; that assumption turns an impossible joint estimate into
simple per-feature counts.}
\why
Start: $P(C\mid\mathbf x) = P(\mathbf x\mid C)P(C)/P(\mathbf x)$. The denominator is the same for all classes, so we compare numerators.
The hard part is $P(\mathbf x\mid C)$ for a high-dimensional $\mathbf x$: with 10{,}000 binary word features there are $2^{10000}$
combinations, and we would need data for each. Conditional independence lets us write it as a product of 10{,}000 one-word
probabilities, each estimated by counting. Note the independence is \emph{conditional on the class}, not overall: ``free'' and
``offer'' are clearly correlated overall because both are spammy; the assumption is that within spam they occur independently.
In practice: Multinomial NB on word counts with Laplace smoothing, computed in log space, as a fast baseline for text.
\worked The ``free meeting'' example: spam score 0.006, ham 0.0105, posterior spam 0.36.
\pushes \emph{``Why does it work when the assumption is false?''} Classification needs only the correct argmax, not correct
probabilities; errors from correlated features often affect all classes similarly. Also, its low variance beats more flexible models
when data are scarce. \emph{``What about the probabilities?''} Over-confident; calibrate if needed.
\pitfall Saying the features must be ``independent'' without ``given the class'', or calling it a discriminative model (it is generative:
it models $P(\mathbf x\mid C)P(C)$).
\end{qa}

\begin{qa}{\pyq{INT: Fisher's LDA} Explain Fisher's linear discriminant: the formulation and the intuition.}
\oneline{Find a projection direction that maximises the squared distance between projected class means divided by the total
within-class variance along that direction; the solution is $\mathbf w\propto S_W^{-1}(\mathbf m_1 - \mathbf m_2)$.}
\why
Projecting onto a line reduces each point to one number. A good line keeps the classes apart (numerator: between-class separation) and
each class compact (denominator: within-class spread). Maximising the ratio $J(\mathbf w) = \frac{\mathbf w^\top S_B\mathbf w}{\mathbf w^\top
S_W\mathbf w}$: since only the direction matters, set the gradient to zero, giving $S_B\mathbf w = \lambda S_W\mathbf w$, a generalised
eigenproblem. Because $S_B\mathbf w$ always points along $\mathbf m_1 - \mathbf m_2$, the answer is $S_W^{-1}(\mathbf m_1 - \mathbf m_2)$. For a
classifier, threshold the projection (e.g. at the midpoint of projected means, adjusted for priors).
\worked See the 2-D example: means difference $(-2,-1)$ becomes Fisher direction $(-0.5,-1)$ after whitening by $S_W^{-1}$.
\pushes \emph{``How is it different from PCA?''} PCA ignores labels and keeps directions of maximum total variance, which may be useless
for classification (the cigars example); LDA uses labels. \emph{``How many components?''} At most $K-1$. \emph{``When does it
fail?''} Non-linear class structure, multimodal classes, very unequal covariances (use QDA), $d > n$ (use shrinkage).
\end{qa}

\begin{qa}{How does $k$-NN work, and how do you choose $k$?}
\oneline{It predicts by majority vote (or average) of the $k$ closest training points; small $k$ gives low bias and high variance, large
$k$ the reverse, so choose $k$ by cross-validation.}
\why Worked example above: 3-NN classifies $Q$ as S. With $k=1$ the decision boundary is the Voronoi tessellation of training points.
\pushes \emph{``Training and prediction cost?''} Training free; prediction $O(nd)$ brute force; use KD-trees or ANN indexes.
\end{qa}

\begin{qa}{\deep Why does $k$-NN suffer from the curse of dimensionality? Give a number.}
\oneline{In high dimensions volumes grow so fast that a neighbourhood holding a fixed fraction of data must span most of each axis, and all
pairwise distances become nearly equal, so ``nearest'' neighbours are not similar.}
\why To contain 1\% of uniform data in $d = 100$ dimensions, a sub-cube needs side $0.01^{1/100} = 0.955$ of the full range. Also the ratio
(farthest $-$ nearest)/nearest distance tends to 0 as $d$ grows for many distributions. Remedies: feature selection, PCA, learned
embeddings (where the meaningful structure is low-dimensional), and metric learning.
\end{qa}

\begin{qa}{Why must features be scaled for $k$-NN, and does it matter for Naive Bayes?}
\oneline{$k$-NN uses distances, so unscaled features with large ranges dominate; Naive Bayes models each feature separately per class,
so scaling does not change its decisions (Gaussian NB just rescales its means and variances).}
\end{qa}

\begin{qa}{Compute the Laplace-smoothed probability of ``offer'' in spam if spam has 20 tokens, vocabulary size 10, and ``offer'' appears
3 times. Why is smoothing needed?}
\oneline{$(3+1)/(20+10) = 0.133$; without smoothing, a word never seen in a class would give probability 0 and wipe out the whole product.}
\end{qa}

\begin{qa}{Why do we compute Naive Bayes in log space?}
\oneline{Products of many small probabilities underflow to 0 in floating point; summing logs is numerically stable and preserves the argmax.}
\worked 200 words each with probability 0.001: product $10^{-600}$, below double precision's $10^{-308}$; log sum $-1381.6$ is fine.
\end{qa}

\begin{qa}{\deep Naive Bayes vs logistic regression: they are both linear for text. What is the difference?}
\oneline{Multinomial Naive Bayes gives a linear decision rule in word counts, just like logistic regression, but NB estimates the weights
generatively from class-conditional counts (fast, low variance, good with little data) while logistic regression fits them
discriminatively to maximise conditional likelihood (better with more data, handles correlated features).}
\why NB's log-odds $=\log\frac{P(S)}{P(H)} + \sum_jx_j\log\frac{P(w_j\mid S)}{P(w_j\mid H)}$: linear in counts $x_j$. Ng and Jordan showed
NB reaches its (higher) asymptotic error faster; logistic regression reaches a lower error with enough data. Generative--discriminative pair.
\end{qa}

\begin{qa}{\deep If two features are exact duplicates, what happens to Naive Bayes? To logistic regression?}
\oneline{Naive Bayes counts the same evidence twice, so its posterior becomes more extreme (over-confident); logistic regression splits the
weight between the duplicates and its predictions are unchanged.}
\why NB multiplies $P(x_1\mid C)P(x_1\mid C)$, squaring the likelihood ratio. Logistic regression's predictions depend on $w_1 + w_2$, which
the fit learns correctly (with L2 it splits evenly). This is the clearest demonstration of what the naive assumption costs.
\end{qa}

\begin{qa}{What is the difference between generative and discriminative classifiers? Classify NB, LDA, logistic regression, SVM, $k$-NN.}
\oneline{Generative models learn how data are produced, $P(\mathbf x\mid y)P(y)$ (Naive Bayes, LDA/QDA, GMM classifiers, GANs and VAEs as
generators); discriminative models learn $P(y\mid\mathbf x)$ or a boundary directly (logistic regression, SVM, trees, neural nets);
$k$-NN is a non-parametric discriminative method.}
\why Generative models can generate samples, handle missing features naturally, and do well with little data; discriminative models usually
win on accuracy with enough data.
\end{qa}

\begin{qa}{LDA vs PCA: when would PCA's first component be useless for classification?}
\oneline{When the direction of largest total variance is unrelated to the class difference, for example two parallel elongated clusters
separated across their narrow direction.}
\why PCA picks the long axis (where both clusters spread); projecting on it mixes the classes completely. LDA picks the narrow axis.
\end{qa}

\begin{qa}{\deep Under what assumptions is LDA the Bayes-optimal classifier, and what changes with QDA?}
\oneline{If each class is Gaussian with the same covariance, the log-posterior ratio is linear in $\mathbf x$ and LDA is optimal; with
different covariances the quadratic terms no longer cancel and the optimal boundary is quadratic (QDA).}
\why $\log\frac{P(1\mid\mathbf x)}{P(2\mid\mathbf x)} = \mathbf x^\top\Sigma^{-1}(\boldsymbol\mu_1 - \boldsymbol\mu_2) + \text{const}$ when the
$\mathbf x^\top\Sigma^{-1}\mathbf x$ terms cancel. QDA needs a covariance per class: $Kd(d+1)/2$ parameters, more variance; regularised DA
interpolates.
\end{qa}

\begin{qa}{How would you speed up nearest-neighbour search over 50 lakh job embeddings?}
\oneline{Use an approximate nearest-neighbour index such as HNSW or IVF with product quantisation (FAISS/ScaNN), trading a little recall for
orders-of-magnitude speed, and pre-filter by hard constraints like location.}
\why HNSW builds a navigable small-world graph; search hops greedily toward the query in $O(\log n)$. IVF clusters vectors and searches only
the nearest clusters; PQ compresses vectors to bytes. Measure recall@k against brute force on a sample.
\end{qa}

\begin{qa}{\deep Can $k$-NN be used for regression and for imputation? What are the risks?}
\oneline{Yes: average the neighbours' targets for regression, and fill a missing value with the neighbours' values for imputation; risks are
scaling, irrelevant features dominating distances, high dimensions, and leakage if the imputer is fitted on test data.}
\worked Neighbours with value 10 at distance 1 and 20 at distance 2, weights $1/d$: $(10 + 10)/(1 + 0.5) = 13.3$.
\end{qa}

\begin{qa}{Gaussian Naive Bayes: how would it classify a candidate with test score 70 if passers have mean 75, sd 5 and failers mean 60,
sd 10, with equal priors?}
\oneline{Compare normal densities: passers $\phi(-1)/5 = 0.0484$, failers $\phi(1)/10 = 0.0242$; predict ``pass'' with posterior $0.667$.}
\why $\phi(1) = 0.242$. Passers: $z = (70-75)/5 = -1$, density $0.242/5 = 0.0484$. Failers: $z = 1$, $0.242/10 = 0.0242$. Ratio 2:1.
\end{qa}

\begin{qa}{\deep Is $k$-NN with $k = 1$ ever a good model?}
\oneline{Yes, when classes are well separated, labels are clean and the distance is meaningful (e.g. near-duplicate detection, or embeddings
trained for similarity); its asymptotic error is at most twice the Bayes error.}
\why The Cover--Hart result: as $n\to\infty$, 1-NN error $\le 2\times$ Bayes error. With good embeddings, 1-NN retrieval is the basis of
duplicate-resume detection and ``find the canonical skill for this raw string'' in taxonomy normalisation.
\end{qa}
